% Prepared protocol text only: no new results or completion claim.
% M1--M8 correspond to the paragraphs in final_program_methods_zh.md.
\subsection{Shared-trajectory checkpoint-selection protocol}
\label{sec:final-program-methods}

% M1: Scope and prior exposure.
The protocol specifies 360 main training trajectories, comprising 24 draws
and five outer folds per corpus, plus 24 paired CREMA-D control fits.
It compares checkpoint-selection rules on a common training trajectory and
common speaker-independent outer test set. CREMA-D, SUBESCO, and RAVDESS
retain their native six-, seven-, and eight-class tasks and 91, 20, and
24 speakers, respectively~\cite{cremad,subesco,ravdess}.
These corpora have prior research exposure;
new draws do not constitute previously untouched data. The protocol uses
a private pre-execution source/plan freeze, not public preregistration.

% M2: Speaker/text allocation and representatives.
Within each draw, metadata-based hashing partitions speakers into five outer
folds, stratified by source-recorded sex; each speaker appears in exactly one
outer fold. Two disjoint development groups, A and B, are selected outside
that fold with equal female/male counts and complete native-class coverage
on the chosen fit and query texts. All three query groups use the same
texts, disjoint from fit texts (Table~\ref{tab:final-program-panels}).
Unselected development speakers remain unused. After the frozen hygiene
rules, CREMA-D uses MD-preferred, otherwise XX representatives; SUBESCO uses
T1; RAVDESS uses normal-intensity R1. Missing predefined representatives
are not replaced by other takes.
Outer queries retain available representatives, requiring every outer
speaker to support all native classes. Insufficient development capacity
causes a metadata failure, not outcome-dependent resampling.

\begin{table}[t]
\centering
\caption{Prespecified per-context allocation. A and B are equal-sized,
sex-balanced groups; all counts refer to distinct speakers or recordings,
not training updates.}
\label{tab:final-program-panels}
\small
\setlength{\tabcolsep}{3pt}
\begin{tabular}{lrrrcrr}
\hline
Corpus & Classes & A/B & Outer & Texts$^a$ & Fit & Query$^b$\\
\hline
CREMA-D & 6 & 24/24 & 17/19 & 4/2 & 576 & 288\\
SUBESCO & 7 & 8/8 & 4 & 4/2 & 224 & 112\\
RAVDESS & 8 & 8/8 & 4/6 & 1/1 & 64 & 64\\
\hline
\end{tabular}
\par\raggedright\small
$^a$ Fit/query text counts. $^b$ Recordings in each of A and B query;
outer queries contain available CREMA-D representatives, 56 SUBESCO
recordings, or 32/48 RAVDESS recordings. Slash-separated outer counts
denote alternative fold sizes.
\end{table}

% M3: Fixed optimization and waveform handling.
Only A supplies training recordings for each main trajectory. WavLM
base-plus~\cite{wavlm} uses a native-class linear head over mean-pooled final-layer
features; only the top four Transformer layers and head are trainable.
The fixed configuration, denoted cfg3 in the earlier study, uses AdamW,
encoder/head learning rates $5\times10^{-5}/10^{-3}$, weight decay 0.01,
batch size 16, and unweighted cross-entropy. Training traverses the fit
set for all 15 epochs without early stopping or hyperparameter search,
using FP16 autocast with FP32 parameters; evaluation uses FP32.
Audio is loaded as mono at 16\,kHz, silence-trimmed at 30\,dB, cached
in FP16 and converted to FP32 for batching. Training uses three-second
crops or zero-padding for shorter clips; evaluation takes up to the first
ten seconds of processed audio.

% M4: Four rules, tie handling, and fixed padding context.
For A-trained models, A queries are seen-speaker validation and B queries
are unseen-speaker validation. Each epoch records logits for A, B, and
outer queries. Four rules minimize mean cross-entropy or maximize
unweighted average recall (UAR), separately on A or B; exact ties select
the earliest epoch. UAR is macro recall across all native classes in the
whole query group, not a mean of speaker UARs. Integer/rational recall
preserves exact ties. Each group has its own fixed sorted batches of 16,
reused across epochs and paired models. Outer recordings never pad
validation batches. Because pooling lacks a length mask, this fixes
padding context without eliminating padding effects. Evaluation preserves
training RNG state, and training mode is restored each epoch.

% M5: Estimands.
Let $T(g,q)$ denote outer UAR, in percent, at the epoch selected by
validation group $g$ (seen $s$ or unseen $u$) and criterion
$q\in\{\mathrm{CE},\mathrm{UAR}\}$.
Four rules reuse one fitted trajectory; they are not four training runs.
The per-context contrasts are
\begin{align}
\delta_{\mathrm{CE}} &= T(s,\mathrm{CE})-T(u,\mathrm{CE}),\\
\delta_{\mathrm{UAR}} &= T(s,\mathrm{UAR})-T(u,\mathrm{UAR}),\\
J &= \delta_{\mathrm{CE}}-\delta_{\mathrm{UAR}}.
\end{align}
These contrasts measure outer UAR differences in percentage points (pp).
They are not differences in cross-entropy loss or in validation optimism
($V-T$).
They compare algorithmic selection rules rather than identify a unique
speaker mechanism.

% M6: Prespecified inferential family and scope.
Within each corpus, contrasts are averaged equally over five folds per
draw. Two-sided one-sample $t$ inference then uses the 24 complete draw
effects ($df=23$). The fixed six-test family is three corpora times
$(\delta_{\mathrm{CE}},J)$, with Holm control at 0.05 and pointwise 95\%
$t$ intervals. Undefined zero-variance tests remain NA and retain their
family slots. Absolute UAR and $\delta_{\mathrm{UAR}}$ are descriptive.
There is no bootstrap or equivalence test; 1\,pp is a prespecified
practical reference, not an equivalence margin. Inference is conditional
on each finite corpus and training program: folds, recordings, repeated
speakers, and rules are not additional replicates. Native tasks differ,
so corpora are neither pooled for inference nor treated as language effects.

% M7: Paired whole-group role control.
For each CREMA-D draw, prespecified fold 0 adds one B-trained model.
Paired fits use identical initial weights and matched seeds in separate
streams for initialization, ordering, cropping, and training; fit slots match sex, within-group speaker
rank, text, class, take, and intensity. B trains for 15 epochs but records
queries only at epoch 15, without checkpoint selection. With $Q_A$ and
$Q_B$ denoting whole-group query UAR in percent, the descriptive control is
\begin{equation}
\begin{aligned}
E=\tfrac12\{&[Q_A(M_{A,15})-Q_A(M_{B,15})]\\
&+[Q_B(M_{B,15})-Q_B(M_{A,15})]\}.
\end{aligned}
\end{equation}
The protocol reports 24 paired values and their mean, without additional
inference. This replaces whole training groups, not individual speakers;
the same queries also participate in main-A checkpoint selection.

% M8: Sealing and prespecified diagnostics.
Four separate technical pilots are outside the 384-fit matrix. Scientific
scoring requires all formal units to pass source, artifact, checkpoint,
and ledger verification. Main trajectories retain all 15 epochs' query
logits and deduplicated rule-winning/final checkpoint deltas; controls
retain the final state. Disk reload requires maximum absolute logit
differences no greater than $10^{-5}$.
Post-sealing descriptions include selected-epoch distributions and
agreement, complete outer curves, last-epoch UAR, epochs 8--10 minus
14--15 mean UAR, and shortfall from each trajectory's maximum outer UAR.
The outer oracle never selects retained checkpoints or extends training.
